% birkon/cover/ornament-options.tex
% Original cover-ornament options for the birkon, one per page, at
% label size (3.25in x 4.5in, for the front board of the 4.25in x 5.5in book).
% For choosing a design (A-F); the chosen one becomes the cover label or wrap.
\documentclass{article}
\IfFileExists{../../shared/macros/hebrew-layout.sty}
  {\def\pressroot{../../}}
  {\def\pressroot{}}
\makeatletter
\def\input@path{{\pressroot shared/macros/}}
\makeatother
\usepackage[paperwidth=3.25in,paperheight=4.5in,margin=0pt]{geometry}
\usepackage{hebrew-layout}
\usepackage{tikz}
\usetikzlibrary{calc}

% Each design is an ordinary tikzpicture exactly the size of the label.
% (Text inside "remember picture, overlay" pictures is lost under babel's
% right-to-left mode, so overlays are not used.)
\def\W{8.255}  % 3.25in in cm
\def\H{11.43}  % 4.5in in cm
% The picture sits in a left-to-right box: in a right-to-left paragraph pgf's
% drawing coordinates are mirrored off the page. Text in nodes is set back
% to right to left.
\newenvironment{coverlabel}{\noindent\hbox dir TLT\bgroup\begin{tikzpicture}[
    every node/.style={execute at begin node={\pardir TRT\textdir TRT}}]
  \useasboundingbox (0,0) rectangle (\W,\H);
  \coordinate (C) at (\W/2,\H/2);}{\end{tikzpicture}\egroup\newpage}
\pagestyle{empty}
\setlength{\parindent}{0pt}
\setlength{\topskip}{0pt}

\newcommand{\coverTitle}{%
  {\fontsize{19}{24}\selectfont\bfseries סֵדֶר}\\[2pt]%
  {\fontsize{19}{24}\selectfont\bfseries בִּרְכַּת הַמָּזוֹן}}
\newcommand{\coverSub}{%
  \foreignlanguage{english}{\small\itshape Birkat HaMazon}}

% One wheat ear: a stalk along a curve and paired grains.
% #1 = mirror (1 or -1)
\newcommand{\wheatEar}[1]{%
  \begin{scope}[xscale=#1]
    \draw[line width=0.7pt] (0,0) .. controls (0.35,0.25) and (0.75,0.75) .. (0.9,1.45);
    \foreach \t/\a in {0.30/38,0.42/45,0.54/52,0.66/58,0.78/63,0.90/68}{
      \path (0,0) .. controls (0.35,0.25) and (0.75,0.75) .. (0.9,1.45)
        coordinate[pos=\t] (g);
      \fill[rotate around={\a+35:(g)}] ($(g)+(0.07,0)$) ellipse (0.075 and 0.03);
      \fill[rotate around={\a-35+180:(g)}] ($(g)+(-0.07,0)$) ellipse (0.075 and 0.03);
    }
    \path (0,0) .. controls (0.35,0.25) and (0.75,0.75) .. (0.9,1.45) coordinate[pos=1] (tip);
    \fill[rotate around={68:(tip)}] (tip) ellipse (0.07 and 0.03);
  \end{scope}}

% Eight-pointed star rosette, radius #1 (cm).
\newcommand{\rosette}[1]{%
  \draw[line width=0.6pt] (45:#1) -- (135:#1) -- (225:#1) -- (315:#1) -- cycle;
  \draw[line width=0.6pt] (0:#1*1.414/1.414) -- (90:#1) -- (180:#1) -- (270:#1) -- cycle;
  \draw[line width=0.5pt] (0,0) circle (#1*0.42);
  \fill (0,0) circle (#1*0.10);}

\def\skylineWidth{7.6}
% Seattle skyline as one continuous line: the outline of the Space Needle
% (#1 = 1; 0 leaves it out), the Great Wheel, downtown towers, Smith Tower's
% pointed cap and Mount Rainier, starting and ending on the ground. Origin at
% the left end; \skylineWidth wide. \skylinePath continues a path from (0,0)
% and ends at (\skylineWidth,0), so a frame can run into it.
\newcommand{\seattleSkyline}[1]{%
  \draw[line width=0.6pt,line join=round,line cap=round]
    (0,0) \skylinePath{#1};}
\newcommand{\skylinePath}[1]{%
    \ifnum#1=1
      -- (1.11,0) -- (1.30,1.45) -- (1.22,2.0) -- (1.05,2.0) -- (0.93,2.08)
      -- (1.05,2.17) -- (1.25,2.17) -- (1.28,2.27) -- (1.35,2.27) -- (1.35,2.62)
      -- (1.35,2.27) -- (1.42,2.27) -- (1.45,2.17) -- (1.65,2.17) -- (1.77,2.08)
      -- (1.65,2.0) -- (1.48,2.0) -- (1.40,1.45) -- (1.59,0)
    \fi
    % the Great Wheel: legs, spokes (out and back from the hub), rim
    -- (1.80,0) -- (2.05,0.46)
    \foreach \a in {0,45,...,315}{ -- ++(\a:0.33) -- (2.05,0.46) }
    -- ++(0:0.33) arc[start angle=0,end angle=360,radius=0.33]
    -- (2.05,0.46) -- (2.30,0)
    % downtown, ending with Smith Tower's point
    -- (2.55,0) -- (2.55,1.35) -- (2.95,1.35) -- (2.95,2.05) -- (3.12,2.25)
    -- (3.3,2.25) -- (3.47,2.05) -- (3.47,1.62) -- (3.85,1.62) -- (3.85,1.05)
    -- (3.96,1.05) -- (4.12,1.175) -- (4.12,1.25) -- (4.24,1.25) -- (4.31,1.55)
    -- (4.354,1.36) -- (4.354,0.70)
    % Rainier: broad and low, long concave slopes, a shoulder on the left,
    % a flattened summit with a small notch
    .. controls (4.8,0.75) and (5.1,0.92) .. (5.35,1.08) -- (5.45,1.10)
    .. controls (5.6,1.2) and (5.75,1.3) .. (5.9,1.36)
    -- (5.98,1.38) -- (6.04,1.34) -- (6.12,1.38) -- (6.22,1.36)
    .. controls (6.5,1.15) and (6.85,0.35) .. (\skylineWidth,0)}

\newcommand{\ornDiamond}[1]{\fill (#1) +(0,0.06) -- +(0.06,0) -- +(0,-0.06) -- +(-0.06,0) -- cycle;}

\begin{document}

% ---- A: single frame, notched corners -------------------------------------
\begin{coverlabel}
  \def\r{0.32}\def\m{1.15}
  \begin{scope}[line width=0.7pt]
  \draw (\m,\m+\r) arc[start angle=90,end angle=0,radius=\r]
     -- (\W-\m-\r,\m) arc[start angle=180,end angle=90,radius=\r]
     -- (\W-\m,\H-\m-\r) arc[start angle=270,end angle=180,radius=\r]
     -- (\m+\r,\H-\m) arc[start angle=0,end angle=-90,radius=\r]
     -- cycle;
  \end{scope}
  \ornDiamond{\W/2,\m}\ornDiamond{\W/2,\H-\m}%
  \ornDiamond{\m,\H/2}\ornDiamond{\W-\m,\H/2}
  \node[align=center] at (C) {\coverTitle\\[6pt]%
    \rule[0.5ex]{1.4cm}{0.5pt}\\[4pt]\coverSub};
  \node[font=\scriptsize] at (\W/2,0.45) {A};
\end{coverlabel}

% ---- B: wheat ----------------------------------------------------------------
\begin{coverlabel}
  \node[align=center] at ($(C)+(0,0.7)$) {\coverTitle\\[6pt]\coverSub};
  \begin{scope}[shift={($(C)+(0,-2.3)$)},scale=1.25]
    \begin{scope}[shift={(0.12,0)}]\wheatEar{1}\end{scope}
    \begin{scope}[shift={(-0.12,0)}]\wheatEar{-1}\end{scope}
  \end{scope}
  \node[font=\scriptsize] at (\W/2,0.45) {B};
\end{coverlabel}

% ---- C: eight-pointed star rosette ---------------------------------------
\begin{coverlabel}
  \begin{scope}[shift={($(C)+(0,2.6)$)}]\rosette{0.75}\end{scope}
  \node[align=center] at ($(C)+(0,0.2)$) {\coverTitle};
  \draw[line width=0.5pt] ($(C)+(-1.2,-1.4)$) -- ($(C)+(1.2,-1.4)$);
  \ornDiamond{$(C)+(-1.32,-1.4)$}
  \ornDiamond{$(C)+(1.32,-1.4)$}
  \node at ($(C)+(0,-1.85)$) {\coverSub};
  \node[font=\scriptsize] at (\W/2,0.45) {C};
\end{coverlabel}

% ---- D: Seattle skyline in one line, with the Space Needle ------------------
\begin{coverlabel}
  \node[align=center] at ($(C)+(0,2.2)$) {\coverTitle\\[6pt]\coverSub};
  \begin{scope}[shift={(\W/2-\skylineWidth/2,1.75)}]\seattleSkyline{1}\end{scope}
  \node[font=\scriptsize] at (\W/2,0.45) {D};
\end{coverlabel}

% ---- E: Seattle skyline in one line, without the Space Needle ---------------
\begin{coverlabel}
  \node[align=center] at ($(C)+(0,2.2)$) {\coverTitle\\[6pt]\coverSub};
  \begin{scope}[shift={(\W/2-\skylineWidth/2,1.75)}]\seattleSkyline{0}\end{scope}
  \node[font=\scriptsize] at (\W/2,0.45) {E};
\end{coverlabel}

% ---- F: frame A whose bottom edge becomes the skyline ---------------------
% One line: the frame's bottom edge is the skyline of D, scaled to fit
% between the lower notches.
\begin{coverlabel}
  \def\r{0.32}\def\m{1.15}
  \pgfmathsetmacro\s{(\W-2*\m-2*\r)/\skylineWidth}
  \begin{scope}[line width=0.7pt,line join=round]
  \draw (\m+\r,\m)
     {[shift={(\m+\r,\m)},scale=\s] \skylinePath{1}}
     arc[start angle=180,end angle=90,radius=\r]
     -- (\W-\m,\H-\m-\r) arc[start angle=270,end angle=180,radius=\r]
     -- (\m+\r,\H-\m) arc[start angle=0,end angle=-90,radius=\r]
     -- (\m,\m+\r) arc[start angle=90,end angle=0,radius=\r]
     -- cycle;
  \end{scope}
  \ornDiamond{\W/2,\H-\m}%
  \ornDiamond{\m,\H/2}\ornDiamond{\W-\m,\H/2}
  \node[align=center] at ($(C)+(0,0.75)$) {\coverTitle\\[6pt]%
    \rule[0.5ex]{1.4cm}{0.5pt}\\[4pt]\coverSub};
  \node[font=\scriptsize] at (\W/2,0.45) {F};
\end{coverlabel}

\end{document}
